% Options for packages loaded elsewhere
\PassOptionsToPackage{unicode}{hyperref}
\PassOptionsToPackage{hyphens}{url}
\PassOptionsToPackage{dvipsnames,svgnames,x11names}{xcolor}
\documentclass[
  10pt,
  fontset=fandol]{ctexart}
\usepackage{xcolor}
\usepackage[margin=16mm]{geometry}
\usepackage{amsmath,amssymb}
\setcounter{secnumdepth}{-\maxdimen} % remove section numbering
\usepackage{iftex}
\ifPDFTeX
  \usepackage[T1]{fontenc}
  \usepackage[utf8]{inputenc}
  \usepackage{textcomp} % provide euro and other symbols
\else % if luatex or xetex
  \usepackage{unicode-math} % this also loads fontspec
  \defaultfontfeatures{Scale=MatchLowercase}
  \defaultfontfeatures[\rmfamily]{Ligatures=TeX,Scale=1}
\fi
\usepackage{lmodern}
\ifPDFTeX\else
  % xetex/luatex font selection
\fi
% Use upquote if available, for straight quotes in verbatim environments
\IfFileExists{upquote.sty}{\usepackage{upquote}}{}
\IfFileExists{microtype.sty}{% use microtype if available
  \usepackage[]{microtype}
  \UseMicrotypeSet[protrusion]{basicmath} % disable protrusion for tt fonts
}{}
\makeatletter
\@ifundefined{KOMAClassName}{% if non-KOMA class
  \IfFileExists{parskip.sty}{%
    \usepackage{parskip}
  }{% else
    \setlength{\parindent}{0pt}
    \setlength{\parskip}{6pt plus 2pt minus 1pt}}
}{% if KOMA class
  \KOMAoptions{parskip=half}}
\makeatother
\usepackage{longtable,booktabs,array}
\newcounter{none} % for unnumbered tables
\usepackage{calc} % for calculating minipage widths
% Correct order of tables after \paragraph or \subparagraph
\usepackage{etoolbox}
\makeatletter
\patchcmd\longtable{\par}{\if@noskipsec\mbox{}\fi\par}{}{}
\makeatother
% Allow footnotes in longtable head/foot
\IfFileExists{footnotehyper.sty}{\usepackage{footnotehyper}}{\usepackage{footnote}}
\makesavenoteenv{longtable}
\setlength{\emergencystretch}{3em} % prevent overfull lines
\providecommand{\tightlist}{%
  \setlength{\itemsep}{0pt}\setlength{\parskip}{0pt}}
\usepackage{amsmath,amssymb,mathtools,needspace}
\xeCJKDeclareCharClass{CJK}{"2032,"0400->"04FF,"2160->"216B,"2460->"2473,"25A0,"25C7,"25CB,"25CF}
\IfFontExistsTF{Arial Unicode MS}{\xeCJKsetup{AutoFallBack=true}\setCJKfallbackfamilyfont{\CJKrmdefault}{Arial Unicode MS}}{}
\setcounter{secnumdepth}{0}
\setlength{\emergencystretch}{2em}
\allowdisplaybreaks
\usepackage{bookmark}
\IfFileExists{xurl.sty}{\usepackage{xurl}}{} % add URL line breaks if available
\urlstyle{same}
\hypersetup{
  pdftitle={Romberg 算法},
  colorlinks=true,
  linkcolor={Maroon},
  filecolor={Maroon},
  citecolor={Blue},
  urlcolor={Blue},
  pdfcreator={LaTeX via pandoc}}

\title{Romberg 算法}
\author{}
\date{}

\begin{document}
\maketitle

\subsection{4.3 Romberg 算法}\label{romberg-ux7b97ux6cd5}

\subsubsection{4.3.1 梯形法的递推化}\label{ux68afux5f62ux6cd5ux7684ux9012ux63a8ux5316}

上一节介绍的复化求积方法对提高精度是行之有效的，但在使用求积公式之前必须给出合适的步长。步长太大，精度难以保证；步长太小，又会导致计算量的增加。而事先给出一个恰当的步长往往是困难的。

实际计算中常常采用变步长的计算方案，即在步长逐次分半（即步长二分）的过程中，反复利用复化求积公式进行计算，直至所求得的积分值满足精度要求为止。

下面，我们在变步长的过程中探讨梯形法的计算规律。

设将求积区间 \([a,b]\) 分成 \(n\) 等份，则一共有 \(n+1\) 个分点，按梯形公式（4.2.9）计算积分值 \(T_n\)，需要提供 \(n+1\) 个函数值。如果将求积区间再二分一次，则分点增至 \(2n+1\) 个。将二分前后两个积分值联系起来加以考察，注意到每个子区间 \([x_k,x_{k+1}]\) 经过二分只增加了一个分点 \(x_{k+\frac12}=\dfrac12(x_k+x_{k+1})\)，用复化梯形公式求得该子区间上的积分值为

\[
\frac h4[f(x_k)+2f(x_{k+\frac12})+f(x_{k+1})].
\]

注意，这里 \(h=\dfrac{b-a}{n}\) 代表二分前的步长。将每个子区间上的积分值相加，得

\[
T_{2n}=\frac h4\sum_{k=0}^{n-1}[f(x_k)+f(x_{k+1})]+\frac h2\sum_{k=0}^{n-1}f(x_{k+\frac12}),
\]

从而利用式（4.2.9）可导出下列递推公式：

\[
T_{2n}=\frac12T_n+\frac h2\sum_{k=0}^{n-1}f(x_{k+\frac12}).
\tag{4.3.1}
\]

\textbf{例 4.2} 计算积分值 \(I=\displaystyle\int_0^1\frac{\sin x}{x}\,\mathrm{d}x\)。

\textbf{解} 先对整个区间 \([0,1]\) 使用梯形公式。对于函数 \(f(x)=\dfrac{\sin x}{x}\)，它在 \(x=0\) 的值定义为 \(f(0)=1\)，而 \(f(1)=0.841\,470\,9\)，据梯形公式计算得

\[
T_1=\frac12[f(0)+f(1)]=0.920\,735\,5.
\]

然后将区间二等分，再求出中点的函数值 \(f(1/2)=0.958\,851\,0\)，从而利用递推公式（4.3.1），有

\href{../../source-images/pdf-102.jpeg}{原页图像}

\[
T_2=\frac12T_1+\frac12f\left(\frac12\right)=0.939\,793\,3.
\]

进一步二分求积区间，并计算新分点上的函数值，得

\[
f(1/4)=0.989\,615\,8,\quad f(3/4)=0.908\,851\,6.
\]

再利用式（4.3.1），有

\[
T_4=\frac12T_2+\frac14\left[f\left(\frac14\right)+f\left(\frac34\right)\right]=0.944\,513\,5.
\]

这样不断二分下去，计算结果如表 4.3 所示（表中 \(k\) 代表二分次数，区间等份数 \(n=2^k\)）。

表 4.3

{\def\LTcaptype{none} % do not increment counter
\begin{longtable}[]{@{}
  >{\raggedright\arraybackslash}p{(\linewidth - 10\tabcolsep) * \real{0.1667}}
  >{\raggedright\arraybackslash}p{(\linewidth - 10\tabcolsep) * \real{0.1667}}
  >{\raggedright\arraybackslash}p{(\linewidth - 10\tabcolsep) * \real{0.1667}}
  >{\raggedright\arraybackslash}p{(\linewidth - 10\tabcolsep) * \real{0.1667}}
  >{\raggedright\arraybackslash}p{(\linewidth - 10\tabcolsep) * \real{0.1667}}
  >{\raggedright\arraybackslash}p{(\linewidth - 10\tabcolsep) * \real{0.1667}}@{}}
\toprule\noalign{}
\begin{minipage}[b]{\linewidth}\raggedright
\(k\)
\end{minipage} & \begin{minipage}[b]{\linewidth}\raggedright
1
\end{minipage} & \begin{minipage}[b]{\linewidth}\raggedright
2
\end{minipage} & \begin{minipage}[b]{\linewidth}\raggedright
3
\end{minipage} & \begin{minipage}[b]{\linewidth}\raggedright
4
\end{minipage} & \begin{minipage}[b]{\linewidth}\raggedright
5
\end{minipage} \\
\midrule\noalign{}
\endhead
\bottomrule\noalign{}
\endlastfoot
\(T_n\) & \(0.939\,793\,3\) & \(0.944\,513\,5\) & \(9.945\,690\,9\) & \(0.945\,985\,0\) & \(0.946\,059\,6\) \\
\end{longtable}
}

{\def\LTcaptype{none} % do not increment counter
\begin{longtable}[]{@{}
  >{\raggedright\arraybackslash}p{(\linewidth - 10\tabcolsep) * \real{0.1667}}
  >{\raggedright\arraybackslash}p{(\linewidth - 10\tabcolsep) * \real{0.1667}}
  >{\raggedright\arraybackslash}p{(\linewidth - 10\tabcolsep) * \real{0.1667}}
  >{\raggedright\arraybackslash}p{(\linewidth - 10\tabcolsep) * \real{0.1667}}
  >{\raggedright\arraybackslash}p{(\linewidth - 10\tabcolsep) * \real{0.1667}}
  >{\raggedright\arraybackslash}p{(\linewidth - 10\tabcolsep) * \real{0.1667}}@{}}
\toprule\noalign{}
\begin{minipage}[b]{\linewidth}\raggedright
\(k\)
\end{minipage} & \begin{minipage}[b]{\linewidth}\raggedright
6
\end{minipage} & \begin{minipage}[b]{\linewidth}\raggedright
7
\end{minipage} & \begin{minipage}[b]{\linewidth}\raggedright
8
\end{minipage} & \begin{minipage}[b]{\linewidth}\raggedright
9
\end{minipage} & \begin{minipage}[b]{\linewidth}\raggedright
10
\end{minipage} \\
\midrule\noalign{}
\endhead
\bottomrule\noalign{}
\endlastfoot
\(T_n\) & \(0.946\,076\,9\) & \(0.946\,081\,5\) & \(0.946\,082\,7\) & \(0.946\,083\,0\) & \(0.946\,083\,1\) \\
\end{longtable}
}

积分 \(I\) 的准确值为 \(0.946\,083\,1\)，用变步长方法二分 9 次得到了这个结果。

\begin{quote}
校注（agent 补充，CH04A-E03）：表 4.3 的 \(k=3\) 单元格原扫描确实印为 \(9.945\,690\,9\)，这里按原表保留。由例 4.1 的 \(T_8\) 和本页下段的 \(T_8\) 可知首位应为 \(0\)，即 \(0.945\,690\,9\)，这是原书排印错误。

校注（agent 补充，CH04A-E04）：本页说明 \(k\) 为二分次数，表中 \(k=9\) 为 \(0.946\,083\,0\)，\(k=10\) 才列出 \(0.946\,083\,1\)。故表后``二分 9 次得到了这个结果''与本页表格、次数定义不一致；原句已保留。

校注（agent 补充，CH04A-E07）：表 4.3 的 \(k=5\) 单元格原扫描印为 \(0.946\,059\,6\)，这里按原表保留。对本例 \(f(x)=\sin x/x\)（\(f(0)=1\)）在 \([0,1]\) 上取 \(n=32\) 独立高精度计算复化梯形值，得 \(T_{32}=0.946\,058\,560\,962\,768\,067\ldots\)，保留七位小数应为 \(0.946\,058\,6\)；原表数值有误。
\end{quote}

\subsubsection{4.3.2 Romberg 公式}\label{romberg-ux516cux5f0f}

梯形法的算法简单，但精度较差，收敛的速度缓慢。如何提高收敛速度以节省计算量，自然是人们极为关心的课题。

根据梯形法的误差公式（4.2.15），积分值 \(T_n\) 的截断误差大致与 \(h^2\) 成正比，因此当步长二分后，截断误差将减至原有误差的 \(1/4\)，即有

\[
\frac{I-T_{2n}}{I-T_n}\approx\frac14.
\]

将上式移项整理，可得

\[
I-T_{2n}\approx\frac13(T_{2n}-T_n).
\tag{4.3.2}
\]

由此可见，只要二分前后的两个积分值 \(T_n\) 与 \(T_{2n}\) 相当接近，就可以保证计算结果 \(T_{2n}\) 的误差很小。这种直接用计算结果来估计误差的方法通常称为误差的事后估计法。

按式（4.3.2），积分近似值 \(T_{2n}\) 的误差大致等于 \(\dfrac13(T_{2n}-T_n)\)，因此，如果用这个误差值作为 \(T_{2n}\) 的一种补偿，可以期望，所得到的

\[
\overline T=T_{2n}+\frac13(T_{2n}-T_n)=\frac43T_{2n}-\frac13T_n
\tag{4.3.3}
\]

可能是更好的结果。

再考察例 4.2，所求得的两个梯形值 \(T_4=0.944\,513\,5\) 和 \(T_8=0.945\,690\,9\) 的精

\href{../../source-images/pdf-103.jpeg}{原页图像}

度都很差（与准确值 \(I=0.946\,083\,1\) 比较，只有两三位有效数字），但如果将它们按式（4.3.3）作线性组合，则新的近似值

\[
\overline T=\frac43T_8-\frac13T_4=0.946\,083\,3
\]

却有六位有效数字。

按公式（4.3.3）组合得到的近似值 \(\overline T\)，其实质究竟是什么呢？直接验证易知

\[
S_n=\frac43T_{2n}-\frac13T_n,
\tag{4.3.4}
\]

这就是说，用梯形法二分前后的两个积分值 \(T_n\) 与 \(T_{2n}\)，按式（4.3.3）作线性组合，结果得到 Simpson 法的积分值 \(S_n\)。

再考察 Simpson 法，按误差公式（4.2.16），其截断误差大致与 \(h^4\) 成正比，因此，若将步长折半，则误差将减至原有误差的 \(1/16\)，即有

\[
\frac{I-S_{2n}}{I-S_n}\approx\frac1{16},\quad I\approx\frac{16}{15}S_{2n}-\frac1{15}S_n.
\]

不难直接验证，上式右端的值其实等于 \(C_n\)，就是说，用 Simpson 法二分前后的两个积分值 \(S_n\) 与 \(S_{2n}\)，按式（4.3.3）作线性组合，结果得到 Cotes 法的积分值 \(C_n\)，即

\[
C_n=\frac{16}{15}S_{2n}-\frac1{15}S_n.
\tag{4.3.5}
\]

重复同样的手续，依据 Cotes 法的误差公式（4.2.17）可进一步导出下列 Romberg 公式：

\[
R_n=\frac{64}{63}C_{2n}-\frac1{63}C_n.
\tag{4.3.6}
\]

在变步长的过程中运用式（4.3.4）、式（4.3.5）和式（4.3.6），就能将粗糙的梯形值 \(T_n\) 逐步加工成精度较高的 Simpson 值 \(S_n\)、Cotes 值 \(C_n\) 和 Romberg 值 \(R_n\)。

\textbf{例 4.3} 用加速公式（4.3.4）、（4.3.5）和（4.3.6）加工例 4.2 得到的梯形值，计算结果如表 4.4（\(k\) 代表二分次数）所示。

表 4.4

{\def\LTcaptype{none} % do not increment counter
\begin{longtable}[]{@{}
  >{\raggedright\arraybackslash}p{(\linewidth - 8\tabcolsep) * \real{0.2000}}
  >{\raggedright\arraybackslash}p{(\linewidth - 8\tabcolsep) * \real{0.2000}}
  >{\raggedright\arraybackslash}p{(\linewidth - 8\tabcolsep) * \real{0.2000}}
  >{\raggedright\arraybackslash}p{(\linewidth - 8\tabcolsep) * \real{0.2000}}
  >{\raggedright\arraybackslash}p{(\linewidth - 8\tabcolsep) * \real{0.2000}}@{}}
\toprule\noalign{}
\begin{minipage}[b]{\linewidth}\raggedright
\(k\)
\end{minipage} & \begin{minipage}[b]{\linewidth}\raggedright
\(T_2^k\)
\end{minipage} & \begin{minipage}[b]{\linewidth}\raggedright
\(S_2^{k-1}\)
\end{minipage} & \begin{minipage}[b]{\linewidth}\raggedright
\(C_2^{k-2}\)
\end{minipage} & \begin{minipage}[b]{\linewidth}\raggedright
\(R_2^{k-3}\)
\end{minipage} \\
\midrule\noalign{}
\endhead
\bottomrule\noalign{}
\endlastfoot
0 & \(0.920\,735\,5\) & & & \\
1 & \(0.939\,793\,3\) & \(0.946\,145\,9\) & & \\
2 & \(0.944\,513\,5\) & \(0.946\,086\,9\) & \(0.946\,083\,0\) & \\
3 & \(0.945\,690\,9\) & \(0.946\,083\,3\) & \(0.946\,083\,1\) & \(0.946\,083\,1\) \\
\end{longtable}
}

我们看到，这里利用二分 3 次的数据（它们的精度都很差，只有两三位是有效数字），通过三次加速求得 \(R_1=0.946\,083\,1\)，这个结果的每一位数字都是有效数字，可见加速效果是十分显著的。

\begin{quote}
校注（agent 补充）：表 4.4 的上、下标按原扫描字位保留为 \(T_2^k\)、\(S_2^{k-1}\)、\(C_2^{k-2}\)、\(R_2^{k-3}\)。由本页``\(k\) 代表二分次数''及公式（4.3.4）至（4.3.6），各列数据对应的分段数实际上为 \(2^k\)、\(2^{k-1}\)、\(2^{k-2}\)、\(2^{k-3}\)，即通常记作 \(T_{2^k}\)、\(S_{2^{k-1}}\)、\(C_{2^{k-2}}\)、\(R_{2^{k-3}}\)。另本页由 Simpson 值组合为 Cotes 值的文字仍引用``式（4.3.3）''，已按原文保留；该处实际所列权为 \(16/15\) 和 \(-1/15\)，见式（4.3.5）。
\end{quote}

\href{../../source-images/pdf-104.jpeg}{原页图像}

\subsubsection{4.3.3 Richardson 外推加速法}\label{richardson-ux5916ux63a8ux52a0ux901fux6cd5}

上述加速过程还可以再继续下去，其理论依据是梯形法的余项可展开成下列级数形式。

\textbf{定理 4.3} 设 \(f(x)\in C^\infty[a,b]\)，则成立

\[
T(h)=I+a_1h^2+a_2h^4+a_3h^6+\cdots+a_kh^{2k}+\cdots,
\tag{4.3.7}
\]

其中系数 \(a_k\)（\(k=1,2,\cdots\)）与 \(h\) 无关。

这一余项公式的推导将在 4.3.4 节给出。

按式（4.3.7），有

\[
T\left(\frac h2\right)=I+\frac{\alpha_1}{4}h^2+\frac{\alpha_2}{16}h^4+\frac{\hat\alpha_3}{64}h^6+\cdots.
\tag{4.3.8}
\]

注意，这里的 \(\hat\alpha_k\) 与将要出现的 \(\beta_k\)，\(\gamma_k\) 等均为与 \(h\) 无关的系数。

设将式（4.3.7）与式（4.3.8）按以下方式作线性组合：

\[
T_1(h)=\frac43T\left(\frac h2\right)-\frac13T(h),
\tag{4.3.9}
\]

则可以从余项展开式中消去误差的主要部分 \(h^2\) 项，从而得到

\[
T_1(h)=I+\beta_1h^4+\beta_2h^6+\beta_3h^8+\cdots.
\tag{4.3.10}
\]

比较式（4.3.9）与式（4.3.4）知，这样构造出的 \(\{T_1(h)\}\) 其实就是 Simpson 值序列。

又根据式（4.3.10），有

\[
T_1\left(\frac h2\right)=I+\frac{\beta_1}{16}h^4+\hat\beta_2h^6+\beta_3h^8+\cdots,
\]

若令

\[
T_2(h)=\frac{16}{15}T_1\left(\frac h2\right)-\frac1{15}T_1(h),
\]

则又可进一步从余项展开式中消去 \(h^4\) 项，从而有

\[
T_2(h)=I+\gamma_1h^6+\gamma_2h^8+\cdots.
\]

这样构造出的 \(\{T_2(h)\}\) 其实就是 Cotes 值序列。

如此继续下去，每加速一次，误差的量级便提高二阶。一般地，将 \(T_0(h)=T(h)\) 按公式

\[
T_m(h)=\frac{4^m}{4^m-1}T_{m-1}\left(\frac h2\right)-\frac1{4^m-1}T_{m-1}(h)
\tag{4.3.11}
\]

经过 \(m\)（\(m=1,2,\cdots\)）次加速后，余项便取下列形式：

\[
T_m(h)=I+\delta_1h^{2(m+1)}+\delta_2h^{2(m+2)}+\cdots.
\tag{4.3.12}
\]

上述处理方法通常称为 Richardson 外推加速方法。

设以 \(T_0^{(k)}\) 表示二分 \(k\) 次后求得的梯形值，且以 \(T_m^{(k)}\) 表示序列 \(\{T_0^{(k)}\}\) 的 \(m\) 次加速值，则依递推公式（4.3.11），即

\[
T_m^{(k)}=\frac{4^m}{4^m-1}T_{m-1}^{(k+1)}-\frac1{4^m-1}T_{m-1}^{(k)}\quad(k=1,2,\cdots)
\tag{4.3.13}
\]

\begin{quote}
校注（agent 补充）：式（4.3.8）的系数按原扫描保留为 \(\alpha_1\)、\(\alpha_2\)、\(\hat\alpha_3\)，相邻正文为 \(\hat\alpha_k\)；\(T_1(h/2)\) 的展开式按原扫描保留 \(\hat\beta_2\) 与未带帽的 \(\beta_3\)。这些字形已放大核对，未用常见写法替换原符号。

校注（agent 补充，CH04A-E05）：式（4.3.13）原列 \(k=1,2,\cdots\)，已照录；但下一页的三角表及算法需要计算 \(T_m^{(0)}\)，因而实现同一递推式时也须允许 \(k=0\)。这是原书递推指标范围的缺漏。

校注（agent 补充，CH04A-E06）：仅有 \(f\in C^\infty[a,b]\) 并不能保证式（4.3.7）的无穷级数收敛且等于 \(T(h)\)。在外推分析中应将其理解为 \(h\to0\) 时的渐近展开，即对所需有限阶截断并保留余项；本转录保留原书等式和条件。下一页式（4.3.14）的无穷 Taylor 等式也有同样的光滑性与收敛性区别。
\end{quote}

\href{../../source-images/pdf-105.jpeg}{原页图像}

可以逐行构造出下列三角形数表------\(T\) 数表：

\[
\begin{array}{ccccc}
T_0^{(0)}&&&&\\
T_0^{(1)}&T_1^{(0)}&&&\\
T_0^{(2)}&T_1^{(1)}&T_2^{(0)}&&\\
T_0^{(3)}&T_1^{(2)}&T_2^{(1)}&T_3^{(0)}&\\
\vdots&\vdots&\vdots&\vdots&\ddots
\end{array}
\]

可以证明，如果 \(f(x)\) 充分光滑，那么 \(T\) 数表每一列的元素及对角线元素均收敛到所求的积分值 \(I\)，即

\[
\lim_{k\to\infty}T_m^{(k)}=I\quad(m\text{ 固定}),\quad\lim_{m\to\infty}T_m^{(0)}=I.
\]

电子计算机上的所谓 Romberg 算法，就是在二分过程中逐步形成 \(T\) 数表的具体方法，其步骤如下：

\textbf{步 1} 准备初值。计算 \(T_0^{(0)}=\dfrac{b-a}{2}[f(a)+f(b)]\) 且令 \(1\to k\)（\(k\) 记录二分的次数）。

\textbf{步 2} 求梯形值。按递推公式（4.3.1）计算梯形值 \(T_0^{(k)}\)。

\textbf{步 3} 求加速值。按加速公式（4.3.13）逐个求出 \(T\) 数表第 \(k+1\) 行其余各元素 \(T_j^{(k-j)}\)（\(j=1,2,\cdots,k\)）。

\textbf{步 4} 精度控制。对于指定精度 \(\varepsilon\)，若 \(|T_k^{(0)}-T_{k-1}^{(0)}|<\varepsilon\)，则终止计算，并取 \(T_k^{(0)}\) 作为所求的结果；否则令 \(k+1\to k\)（意即二分一次），转步 2 继续计算。

\subsubsection{4.3.4 梯形法的余项展开式}\label{ux68afux5f62ux6cd5ux7684ux4f59ux9879ux5c55ux5f00ux5f0f}

前述外推方法的基础是梯形法的余项展开式（4.3.7），现在运用 Taylor 展开方法推导这个公式。

将 \(f(x)\) 在子区间 \([x_k,x_{k+1}]\) 的中点 \(x_{k+\frac12}=\dfrac12(x_k+x_{k+1})\) 展开，有

\[
\begin{aligned}
f(x)={}&f_{k+\frac12}+(x-x_{k+\frac12})f'_{k+\frac12}
+\frac{(x-x_{k+\frac12})^2}{2!}f''_{k+\frac12}
+\frac{(x-x_{k+\frac12})^3}{3!}f'''_{k+\frac12}\\
&+\frac{(x-x_{k+\frac12})^4}{4!}f^{(4)}_{k+\frac12}
+\frac{(x-x_{k+\frac12})^5}{5!}f^{(5)}_{k+\frac12}
+\frac{(x-x_{k+\frac12})^6}{6!}f^{(6)}_{k+\frac12}+\cdots,
\end{aligned}
\tag{4.3.14}
\]

其中 \(f^{(j)}_{k+\frac12}\) 是 \(f^{(j)}(x_{k+\frac12})\) 的缩写。据此写出 \(f(x_k)\) 和 \(f(x_{k+1})\) 的展开式，易知

\[
\frac h2[f(x_k)+f(x_{k+1})]
=hf_{k+\frac12}+\frac h{2!}\left(\frac h2\right)^2f''_{k+\frac12}
+\frac h{4!}\left(\frac h2\right)^4f^{(4)}_{k+\frac12}
+\frac h{6!}\left(\frac h2\right)^6f^{(6)}_{k+\frac12}+\cdots,
\]

求和得

\[
\begin{aligned}
T(h)&=\frac h2\sum[f(x_k)+f(x_{k+1})]\\
&=h\sum f_{k+\frac12}+\frac{h^3}{2!\times2^2}\sum f''_{k+\frac12}
+\frac{h^5}{4!\times2^4}\sum f^{(4)}_{k+\frac12}
+\frac{h^7}{6!\times2^6}\sum f^{(6)}_{k+\frac12}+\cdots.
\end{aligned}
\tag{4.3.15}
\]

\href{../../source-images/pdf-106.jpeg}{原页图像}

\subsubsection{4.3.4 梯形法的余项展开式（续）}\label{ux68afux5f62ux6cd5ux7684ux4f59ux9879ux5c55ux5f00ux5f0fux7eed}

为简洁起见，这里省略了符号 \(\displaystyle\sum_{k=0}^{n-1}\) 中的上、下限。

另一方面，将式（4.3.14）在子区间 \([x_k,x_{k+1}]\) 上求积分，有

\[
\int_{x_k}^{x_{k+1}}f(x)\,\mathrm{d}x
=hf_{k+\frac12}
+\frac{2}{3!}\left(\frac h2\right)^3f''_{k+\frac12}
+\frac{2}{5!}\left(\frac h2\right)^5f^{(4)}_{k+\frac12}
+\frac{2}{7!}\left(\frac h2\right)^7f^{(6)}_{k+\frac12}+\cdots,
\]

然后再关于 \(k\) 从 \(0\) 到 \(n-1\) 求和，得

\[
\begin{aligned}
I&=\int_a^b f(x)\,\mathrm{d}x\\
&=h\sum f_{k+\frac12}
+\frac{h^3}{3!\times2^2}\sum f''_{k+\frac12}
+\frac{h^5}{5!\times2^4}\sum f^{(4)}_{k+\frac12}
+\frac{h^7}{7!\times2^6}\sum f^{(6)}_{k+\frac12}+\cdots.
\end{aligned}
\tag{4.3.16}
\]

利用式（4.3.16）从式（4.3.15）中消去项 \(h\sum f_{k+\frac12}\)，得

\[
T(h)=I+\frac{h^2}{2!\times6}\sum f''_{k+\frac12}
+\frac{h^5}{4!\times20}\sum f^{(4)}_{k+\frac12}
+\frac{3h^7}{6!\times224}\sum f^{(6)}_{k+\frac12}+\cdots.
\tag{4.3.17}
\]

又对 \(f''(x)\) 应用式（4.3.16），并注意到

\[
\int_a^b f''(x)\,\mathrm{d}x=f'(b)-f'(a),
\]

有

\[
h\sum f''_{k+\frac12}=f'(b)-f'(a)
-\frac{h^3}{3!\times2^2}\sum f^{(4)}_{k+\frac12}
-\frac{h^5}{5!\times2^4}\sum f^{(6)}_{k+\frac12}+\cdots,
\]

代入式（4.3.17），整理得

\[
T(h)=I+\frac{h^2}{2!\times6}[f'(b)-f'(a)]
-\frac{h^5}{4!\times30}\sum f^{(4)}_{k+\frac12}
-\frac{h^7}{6!\times56}\sum f^{(6)}_{k+\frac12}+\cdots.
\]

再对 \(f^{(4)}(x)\) 应用式（4.3.16），有

\[
h\sum f^{(4)}_{k+\frac12}=f'''(b)-f'''(a)
-\frac{h^3}{3!\times2^2}\sum f^{(6)}_{k+\frac12}+\cdots,
\]

从而可进一步得

\[
\begin{aligned}
T(h)={}&I+\frac{h^2}{2!\times6}[f'(b)-f'(a)]
-\frac{h^4}{4!\times30}[f'''(b)-f'''(a)]\\
&+\frac{h^7}{6!\times42}\sum f^{(6)}_{k+\frac12}-\cdots.
\end{aligned}
\tag{4.3.18}
\]

反复施行上述手续，即可得到形如式（4.3.7）的余项公式。

最后再强调一下，应用 Romberg 方法时，一定要注意余项展开式（4.3.7）成立这个前提，不然就得不出正确的结果。

\begin{center}\rule{0.5\linewidth}{0.5pt}\end{center}

\subsection{相关原页校注（agent 补充）}\label{ux76f8ux5173ux539fux9875ux6821ux6ce8agent-ux8865ux5145}

以下校注来自本节覆盖的原页，可能也涉及同页相邻小节；原文照录与校正说明分开保留。

\subsubsection{原书 PDF 页 106 的校注}\label{ux539fux4e66-pdf-ux9875-106-ux7684ux6821ux6ce8}

\href{../pages/pdf-106.md}{完整原页转录} · \href{../../source-images/pdf-106.jpeg}{原页图像}

校注记录：CH04B-E001。

\begin{quote}
校注（agent 补充，CH04B-E001）：原扫描式（4.3.17）首个分子明确印为 \(h^2\)，本转录原样保留；该处疑为原书排印错误。由本页式（4.3.16）及中点展开的梯形项相减，\(f''\) 项系数应为 \(h^3/12=h^3/(2!\times6)\)，才能在代入紧接着的 \(h\sum f''\) 关系后得到原页下一式的 \(h^2[f'(b)-f'(a)]/12\)。这不是 OCR 改写；\href{../assets/ch04b/review-p106-middle.png}{局部原图证据}。
\end{quote}

\subsection{本节覆盖原页的校注记录}\label{ux672cux8282ux8986ux76d6ux539fux9875ux7684ux6821ux6ce8ux8bb0ux5f55}

下列记录按原页关联，含同页相邻小节的校注；计算时同时读取上文校注和完整原页。

\begin{itemize}
\tightlist
\item
  \texttt{CH04A-E03}：原书 PDF 页 102
\item
  \texttt{CH04A-E04}：原书 PDF 页 102
\item
  \texttt{CH04A-E05}：原书 PDF 页 104
\item
  \texttt{CH04A-E06}：原书 PDF 页 104
\item
  \texttt{CH04A-E07}：原书 PDF 页 102
\item
  \texttt{CH04B-E001}：原书 PDF 页 106
\end{itemize}

\end{document}
